% ========== 2.5 NON-FUNCTIONAL REQUIREMENTS ==========
% Quality attributes and constraints

\section{2.5 Non-Functional Requirements}
\addcontentsline{toc}{section}{2.5 Non-Functional Requirements}
\label{sec:non-functional-requirements}

\lettrine{N}{on-functional requirements} define quality attributes and constraints that Symphony must satisfy. These requirements specify how well the system performs rather than what it does, addressing performance, reliability, security, and usability characteristics essential for professional development environments.

% ========== NFR Category 1: Performance ==========
\subsection*{2.5.1 Performance}
\addcontentsline{toc}{subsection}{2.5.1 Performance}

\textbf{NFR-1.1 Startup Time:}
System shall achieve rapid startup time on standard hardware configurations, enabling developers to begin work without significant delays.

\textbf{NFR-1.2 Response Latency:}
User interface operations shall respond quickly to maintain fluid interaction experience and developer productivity.

\textbf{NFR-1.3 Infrastructure Operations:}
The Pit extensions shall maintain low operation latency for critical infrastructure functions through Rust implementation.

\textbf{NFR-1.4 Workflow Throughput:}
System shall handle high-volume workflow executions efficiently under normal load conditions.

\textbf{NFR-1.5 Memory Footprint:}
Core system shall maintain minimal memory consumption, with extensions loading on-demand to optimize resource usage.

% ========== NFR Category 2: Reliability & Availability ==========
\subsection*{2.5.2 Reliability \& Availability}
\addcontentsline{toc}{subsection}{2.5.2 Reliability \& Availability}

\textbf{NFR-2.1 System Uptime:}
System shall maintain high availability during normal operation, excluding planned maintenance windows.

\textbf{NFR-2.2 Fault Isolation:}
Extension failures shall not cascade to system-wide crashes through microkernel architecture isolation.

\textbf{NFR-2.3 Data Protection:}
System shall implement automatic saving mechanisms preventing loss of unsaved work during failures.

\textbf{NFR-2.4 Recovery Time:}
System shall recover quickly from component failures through automatic restart and checkpoint resumption mechanisms.

% ========== NFR Category 3: Scalability ==========
\subsection*{2.5.3 Scalability}
\addcontentsline{toc}{subsection}{2.5.3 Scalability}

\textbf{NFR-3.1 Extension Scalability:}
System shall support extensive extension installation without performance degradation through lazy loading mechanisms.

\textbf{NFR-3.2 Project Size:}
System shall efficiently handle large-scale projects with numerous files and substantial codebases.

\textbf{NFR-3.3 Concurrent Workflows:}
System shall support multiple concurrent workflow executions with independent resource allocation.

% ========== NFR Category 4: Security ==========
\subsection*{2.5.4 Security}
\addcontentsline{toc}{subsection}{2.5.4 Security}

\textbf{NFR-4.1 Extension Sandboxing:}
User-facing extensions shall execute in isolated processes with restricted system access.

\textbf{NFR-4.2 Permission System:}
Extensions shall require explicit user approval for accessing sensitive resources or capabilities.

\textbf{NFR-4.3 Audit Logging:}
System shall maintain complete audit trails of extension activities and security-relevant operations.

\textbf{NFR-4.4 Data Privacy:}
System shall provide options for local-only processing without transmitting code to external services.

% ========== NFR Category 5: Usability ==========
\subsection*{2.5.5 Usability}
\addcontentsline{toc}{subsection}{2.5.5 Usability}

\textbf{NFR-5.1 Learning Curve:}
Basic functionality shall be accessible to new users quickly, with intuitive interface design and clear onboarding.

\textbf{NFR-5.2 Interface Consistency:}
User interface shall follow established IDE conventions and maintain consistent interaction patterns.

\textbf{NFR-5.3 Error Messages:}
System shall provide clear, actionable error messages with recovery suggestions.

\textbf{NFR-5.4 Documentation:}
Comprehensive documentation shall be accessible within the system and online.

% ========== NFR Category 6: Compatibility ==========
\subsection*{2.5.6 Compatibility}
\addcontentsline{toc}{subsection}{2.5.6 Compatibility}

\textbf{NFR-6.1 Cross-Platform Support:}
System shall operate on Windows, macOS, and Linux with consistent functionality.

\textbf{NFR-6.2 Language Support:}
System shall support multiple programming languages through extensible syntax highlighting and language servers.


% ========== NFR Category 7: Maintainability ==========
\subsection*{2.5.7 Maintainability}
\addcontentsline{toc}{subsection}{2.5.7 Maintainability}

\textbf{NFR-7.1 Modular Architecture:}
System shall maintain clear separation between core and extensions enabling independent updates.

\textbf{NFR-7.2 Code Quality:}
Codebase shall maintain high quality standards with comprehensive testing and documentation.

\textbf{NFR-7.3 Extension Updates:}
Extensions shall support hot-reloading without system restart for seamless updates.
